\documentclass[11pt,a4paper]{article}
\usepackage[margin=1in]{geometry}
\usepackage{amsmath,amssymb,amsfonts}
\usepackage{mathtools}
\usepackage{lmodern}
\usepackage[T1]{fontenc}
\usepackage[utf8]{inputenc}
\usepackage{textcomp}
\usepackage{microtype}
\usepackage{parskip}
\usepackage{enumitem}
\usepackage{array}
\usepackage{booktabs}
\usepackage{hyperref}
\hypersetup{hidelinks,
  pdftitle={From Primitives to the Substrate-Unruh Effect},
  pdfauthor={Allen Proxmire}}
\setlength{\parindent}{0pt}
\setlength{\parskip}{6pt plus 2pt minus 1pt}

\title{\vspace{-1cm}\Large\bfseries From Primitives to the Substrate-Unruh Effect\\[6pt]
\large\mdseries A Walkthrough of the Event Density Substrate-Level Derivation}
\author{Allen Proxmire}
\date{May 2026}

\begin{document}
\maketitle

\section{The Question}

In 1976, William Unruh published a result that connected accelerated observers to thermal radiation in a way that had no classical analog. A uniformly accelerated observer in flat Minkowski spacetime, moving with proper acceleration $a$, perceives the Minkowski vacuum as a thermal bath at temperature
\[
T_U = a / (2\pi)
\]
(in natural units; with explicit factors, $T_U = \hbar a / (2\pi c k_B)$). At everyday accelerations the temperature is absurdly low --- $a = 9.8$ m/s$^2$ gives $T_U \sim 10^{-19}$ K --- but the structural fact is striking. Vacuum modes that look empty to an inertial observer look thermal to an accelerating observer. The thermalisation has no source: it is not radiation from any external system, it is not noise, it is not a Doppler shift. It is what the vacuum \emph{is} from the accelerated frame's perspective.

The Unruh effect was rapidly recognised as a structurally clean cousin of Hawking radiation. The derivations are formally parallel: in both cases, a Bogoliubov transformation between two natural mode bases (inertial vs.\ accelerated for Unruh, past-vs-future for Hawking) produces a thermal-mixing parameter set by the local acceleration scale (proper acceleration for Unruh, surface gravity for Hawking). The Unruh effect is what the Hawking calculation reduces to in the flat-space limit. The two phenomena share a common structural origin.

The standard derivation runs through QFT in flat spacetime: take the Minkowski vacuum, expand it in modes appropriate to the accelerated observer's frame (Rindler modes), find that the inertial vacuum looks like a thermal state in the accelerated frame at temperature $T_U = a/(2\pi)$. The result is mathematically clean and has been validated indirectly through analog systems (BEC analog accelerating mirrors, accelerating photon detectors) where the thermalisation signature appears at the predicted scale.

But the standard derivation has structural questions that parallel those of standard semiclassical Hawking. The trans-Planckian problem: accelerated modes in the Rindler wedge trace back to arbitrarily-blueshifted near-horizon modes, where standard QFT becomes structurally questionable. The vacuum-state-choice issue: which vacuum (Minkowski, Rindler, Boulware analog) is ``the'' vacuum, and why? The pair-creation interpretation: virtual pairs near the Rindler horizon, with one captured and one observed --- the same ontological heaviness that Hawking derivations carry.

The question this document addresses is: where does the Unruh thermalisation come from at the substrate level, and why does it take exactly the form $T_U = a/(2\pi)$?

The Event Density framework provides a substrate-level account that mirrors the H-arc Hawking derivation. The substrate's V5 vacuum, viewed from a substrate observer accelerated relative to the asymptotic substrate frame at proper substrate-acceleration $a$, exhibits a thermal-correlation structure at temperature $T = a/(2\pi)$. The mechanism is a substrate-level analog of the Euclidean periodicity argument: under analytic continuation of substrate-time, the accelerated substrate frame's geometry has a conical-point-avoidance condition that forces imaginary-time periodicity at $\beta = 2\pi/a$. The KMS condition then identifies this periodicity with thermality at $T = a/(2\pi)$. DCGT identification at leading order maps the substrate calculation to the standard Unruh result; first-subleading corrections from V5 finite-memory cut off the spectrum at the Planck scale.

The chain has six structural moves:

\begin{enumerate}[noitemsep]
\item The substrate has a primitive temporal structure (P13 continuous time, T18 forward-cone-only V1 kernel) and a substrate-level vacuum (the V5 kernel's groundstate).
\item A substrate observer accelerated relative to the asymptotic substrate frame at proper substrate-acceleration $a$ has a substrate-time worldline that traces a hyperbolic trajectory in the substrate's relational structure.
\item The accelerated substrate frame's local geometry, in $(\rho, t)$-like substrate coordinates with $\rho = 1/a$, has a Rindler-like form. Substrate observers at fixed $\rho$ experience proper substrate-acceleration $1/\rho$.
\item Under analytic continuation $t \to -i\tau$, the substrate's local geometry becomes a 2D Euclidean space in polar coordinates, with $a\tau$ as the angular coordinate. Avoiding a substrate-level conical singularity at $\rho = 0$ (the Rindler-like horizon) forces the angular coordinate to have period $2\pi$, giving $\beta = 2\pi/a$.
\item By the KMS condition, V5 correlations periodic in imaginary substrate-time with period $\beta$ are thermal at temperature $T = 1/\beta = a/(2\pi)$.
\item DCGT identification at leading order maps the substrate-level temperature to the standard Unruh temperature exactly. First-subleading corrections from V5 finite-memory at $\tau_{V5} = \ell_P/c$ cut off the high-frequency tail at the Planck scale, providing a substrate-level resolution of the trans-Unruh problem.
\end{enumerate}

The structural payoff: the Unruh effect is what the substrate's V5 vacuum looks like from any accelerated frame. The thermalisation is not a calculational artifact, not a Bogoliubov-mixing technicality, and not a vacuum-state convention. It is the substrate's natural structure expressed through the analytic-continuation periodicity that any accelerated frame imposes on substrate-level correlations.

This walkthrough is the standalone substrate-Unruh derivation. It is structurally identical to the H-1 piece of the Hawking arc, but separable: the substrate-Unruh effect is a substrate-level result on its own, applicable wherever an accelerated frame appears in the substrate (including cosmological horizons, BEC analog systems, and laboratory-accelerated detectors). The Hawking-arc derivation specialises this argument to the saturated-decoupling-surface case; future arcs on de Sitter thermality or other accelerated-frame phenomena will inherit this walkthrough's substrate-Unruh content directly.

\section{The Primitives That Matter}

The framework rests on substrate-level ontological commitments. The substrate-Unruh walkthrough uses the same working subset that the Hawking and Born walkthroughs used.

\textbf{Micro-events (P01).} Discrete acts of becoming.

\textbf{Chains (P02).} Stable subgraphs along which a chain repeatedly instantiates its update rule.

\textbf{Bandwidth (P04).} Non-negative real edge weight with bandwidth-additivity for independent contributions.

\textbf{Polarity / $U(1)$ phase (P09).} $U(1)$-valued phase relation between a chain's update rule and the local ED-flow direction.

\textbf{Continuous time (P13).} The substrate's temporal evolution is continuous between commitment events.

\textbf{Substrate locality.} Cross-chain interactions occur via mediating substrate structure (V1, V5, channels), not via instantaneous non-local action.

\textbf{V1 forward-cone-only kernel (T18).} The substrate's vacuum kernel mediating cross-chain correlations. Forward-cone-only by closed-arc inheritance from the kernel-arrow / arrow-of-time program.

\textbf{V5 finite-memory kernel.} The substrate's cross-chain memory kernel. Has finite memory time $\tau_{V5}$ --- a primitive substrate parameter. V5 supplies the substrate-level vacuum on which the Unruh argument operates.

\textbf{Commitment irreversibility (P11).} Substrate participation events are irreversible.

Two forced theorems load-bear here:

\textbf{T18 (V1 forward-cone-only).} Establishes the substrate's causal structure. Cross-chain correlations propagate forward in substrate-time only.

\textbf{T19 (Newton-recovery $\ell_P$).} Identifies the substrate's irreducible length scale as the Planck length. Load-bears for the V5 timescale identification at the gravitational scale: $\tau_{V5} = \ell_P/c$.

The Diffusion Coarse-Graining Theorem (DCGT) is a structural prerequisite: the substrate-to-continuum bridge that allows the substrate-level temperature to be identified with the standard Unruh temperature at leading order.

That's the structural setup. The substrate-Unruh argument runs on this.

\section{The Substrate Reading of Acceleration}

Standard physics treats acceleration as a kinematical concept: the rate of change of velocity, measured by an inertial observer or experienced by an accelerating one. In ED, acceleration has a substrate-level reading.

\subsection{Substrate-level proper acceleration}

A chain at substrate-time $t$ propagates along a sequence of micro-events with specific spatial and temporal increments. The chain's local \emph{propagation rate} is the substrate quantity governing how rapidly its participation rule advances per unit substrate-time.

When the propagation rate is constant in substrate-time, the chain's worldline traces a straight line in the substrate's emergent spacetime (geodesic at the substrate level). The chain is \emph{inertial}.

When the propagation rate changes with substrate-time, the chain's worldline curves. The chain is \emph{accelerated}. The substrate-level proper acceleration $a$ is the rate of change of the propagation rate per unit substrate-time, evaluated in the chain's instantaneous rest frame.

\subsection{Substrate observers and their accessible regions}

A \emph{substrate observer} is a chain (or chain-aggregate) whose participation rule is sufficiently coherent to integrate substrate events from surrounding regions. The observer's \emph{accessible region} is the substrate region from which V1-mediated cross-chain correlations can reach the observer in finite substrate-time.

For an inertial substrate observer, the accessible region is isotropic --- the V1 forward-cone extends symmetrically around the observer's worldline. The observer integrates substrate events from all directions equally.

For an accelerated substrate observer, the accessible region is \emph{anisotropic}. The observer's V1 forward-cone is tilted along the acceleration direction. Substrate regions in the ``backward'' direction (opposite the acceleration) are progressively excluded; at proper distance $1/a$ behind the observer, the substrate region becomes formally inaccessible --- V1-mediated correlations from that region cannot reach the observer in any finite substrate-time.

\subsection{The Rindler-like horizon}

The proper-distance-$1/a$ surface behind the observer is the substrate analog of the Rindler horizon. It is not a feature of underlying spacetime geometry; it is a feature of the accelerated observer's \emph{accessibility structure}. Different observers in different motion states have different Rindler-like horizons; the surface is observer-dependent.

In the substrate ontology, the Rindler-like horizon is a \emph{participation-bandwidth boundary}: substrate events on the far side cannot integrate into the observer's participation history because the V1 forward-cone, tilted by the observer's acceleration, excludes them.

This is structurally similar to (but distinct from) the saturated decoupling surface of a black-hole horizon. The BH horizon is observer-independent (the gradient is what it is regardless of observer); the Rindler-like horizon is observer-dependent (it appears only in the accelerated frame). Both are participation-bandwidth boundaries; they differ in whether the boundary is set by absolute substrate gradient (BH) or by relative observer-frame structure (Rindler-like).

\subsection{What this delivers}

The substrate-level reading of acceleration produces a Rindler-like substrate frame near the accelerated observer's accessibility horizon. The frame has its own substrate-time direction (the observer's proper substrate-time) and its own substrate-distance coordinate (the proper distance from the horizon).

This Rindler-like substrate frame is the structure on which the substrate-Unruh argument operates. The next sections build the argument on this frame.

\section{The Substrate-Level Vacuum and the Rindler-like Coordinate System}

\subsection{The V5 vacuum}

The framework's V5 cross-chain memory kernel mediates substrate-level correlations between chain regions. The kernel's \emph{vacuum state} is the substrate configuration in which V5 correlations are at their primitive structural level --- the substrate analog of the Wightman vacuum in standard QFT.

Per T18 (forward-cone-only V1), the V5 vacuum has a specific causal structure: cross-chain correlations propagate causally forward in substrate-time. Per the V5 finite-memory primitive, the kernel has characteristic correlation time $\tau_{V5}$ at the gravitational scale (identified as the Planck time $\ell_P/c$ via T19 + dimensional analysis).

The V5 vacuum is not a ``no-particle'' state in the standard QFT sense; it is the substrate's primitive cross-chain correlational structure on which all chain dynamics operate. Substrate observers see this vacuum as the background in which their participation rules unfold.

\subsection{The Rindler-like substrate coordinates}

For a substrate observer accelerated with proper acceleration $a$, introduce substrate coordinates $(\rho, t)$ adapted to the observer's frame:

\begin{itemize}[noitemsep]
\item $\rho$: proper substrate-distance from the observer's Rindler-like horizon at $\rho = 0$.
\item $t$: the observer's proper substrate-time, parameterising the observer's worldline.
\end{itemize}

Substrate observers at fixed $\rho$ experience proper substrate-acceleration $1/\rho$. The originally-considered observer at proper acceleration $a$ sits at $\rho = 1/a$.

\subsection{The local substrate geometry}

Near the Rindler-like horizon at $\rho = 0$, the substrate's local geometry in $(\rho, t)$ coordinates has the form
\[
ds^{2}_{\mathrm{substrate}} \approx -\rho^{2} a^{2} \cdot dt^{2} + d\rho^{2} + \text{(transverse pieces)}
\]
This is the substrate analog of the standard Rindler metric. Substrate observers at fixed $\rho$ experience proper substrate-time elapsed at rate $\rho a$ relative to the asymptotic substrate frame. As $\rho \to 0$, the proper-time rate approaches zero --- the substrate's ``redshift'' effect at the Rindler-like horizon.

\subsection{The asymptotic substrate frame}

The asymptotic substrate frame is the inertial substrate frame at $\rho \to \infty$. In this frame, substrate evolution proceeds at the substrate's natural rate, with no acceleration-induced anisotropy. The V5 vacuum in this frame is the substrate's primitive structural vacuum.

The asymmetry between accelerated and asymptotic frames is structural: the Rindler-like substrate coordinates only cover a wedge of the full substrate (the ``Rindler wedge''), not the entire substrate. Substrate regions outside this wedge are inaccessible to the accelerated observer.

This is the substrate-level analog of the standard Rindler-wedge restriction in QFT. The accelerated observer perceives a \emph{restricted} version of the substrate's V5 vacuum --- restricted to the wedge accessible from the observer's frame.

\section{Imaginary-Time Periodicity (the Load-Bearing Step)}

The substrate-Unruh argument's load-bearing step parallels the standard Euclidean derivation of the Unruh temperature.

\subsection{The Euclidean continuation in substrate-time}

Consider the V5 correlation function across the Rindler-like substrate frame:
\[
G_{V5}(t_1, \rho_1; t_2, \rho_2) \equiv \langle V_5(t_1, \rho_1) V_5^\dagger(t_2, \rho_2)\rangle_{\mathrm{substrate}}
\]
evaluated in the substrate's V5 vacuum. Under analytic continuation $t \to -i\tau$, the correlation function becomes:
\[
G_{V5}^E(\tau_1, \rho_1; \tau_2, \rho_2) \equiv G_{V5}(-i\tau_1, \rho_1; -i\tau_2, \rho_2)
\]
For substrate-states with the appropriate analytic structure (positive-frequency vacuum states, by the substrate analog of the standard QFT positivity-of-spectrum condition), the Euclidean correlation function is well-defined in some strip in the complex plane.

\subsection{The Euclidean substrate geometry}

Under $t \to -i\tau$, the substrate's local geometry near the Rindler-like horizon becomes:
\[
ds^{2}_{\mathrm{substrate}} \approx \rho^{2} a^{2} \cdot d\tau^{2} + d\rho^{2} + \text{(transverse pieces)}
\]
This is the substrate analog of flat 2D Euclidean space in polar coordinates, with $\rho$ as the radial coordinate and $a\tau$ as the angular coordinate.

\subsection{The conical-singularity-avoidance argument}

In standard 2D Euclidean polar coordinates, the metric $ds^{2} = dr^{2} + r^{2} d\theta^{2}$ is regular at $r = 0$ if and only if $\theta$ has period $2\pi$. Otherwise the origin has a conical singularity.

The substrate analog: the substrate's local geometry near $\rho = 0$ is regular only if the angular coordinate $a\tau$ has period $2\pi$. Substrate-time, analytically continued to imaginary substrate-time, must have period
\[
\beta = 2\pi / a
\]
to avoid a substrate-level conical singularity at the Rindler-like horizon.

\subsection{What ``no substrate-level conical singularity'' means}

In standard QFT, the no-conical-singularity argument is geometric: the Euclidean continuation of the Rindler metric must be smooth at the bifurcation surface, which constrains the periodicity of the imaginary time.

The substrate-level statement: V5 cross-chain correlations near the Rindler-like horizon must form a self-consistent substrate object. A conical singularity in the analytically continued substrate would correspond to a discontinuity in V5 correlations as the imaginary substrate-time circle is traversed. Such a discontinuity would mean the substrate's V5 vacuum is not a single coherent substrate object but rather two separate substrate states glued together inconsistently. Substrate locality (P02) and continuity-under-DCGT-coarse-graining forbid such discontinuities at the substrate-vacuum level.

The periodicity $\beta = 2\pi/a$ is therefore \emph{forced} by substrate-vacuum self-consistency at the Rindler-like horizon, in parallel with the standard no-conical-singularity argument in QFT.

\subsection{Status of the periodicity}

The imaginary-time periodicity is FORCED at the substrate level by:
\begin{itemize}[noitemsep]
\item V5 vacuum self-consistency near the Rindler-like horizon
\item Substrate locality and continuity-under-DCGT preventing discontinuities at the substrate-vacuum level
\item The substrate-time invariance of the accelerated frame (the observer's proper time is well-defined and continuous)
\end{itemize}

The argument is structurally parallel to the standard Euclidean-trick derivation of the Unruh temperature, with the substrate's V5 vacuum playing the role of the Wightman vacuum in standard QFT.

\section{The KMS Condition and the Thermal Spectrum}

\subsection{The KMS condition}

A correlation function periodic in imaginary time with period $\beta$ satisfies the \emph{Kubo-Martin-Schwinger} (KMS) condition:
\[
G(-i\beta + \Delta t) = G(\Delta t)^* \quad \text{(with appropriate hermitian conjugation)}
\]
The KMS condition is mathematically equivalent to the correlation function being thermal at temperature $T = 1/\beta$. This is a theorem of standard mathematical physics (Kubo 1957; Martin-Schwinger 1959), inherited as identification target.

The substrate's V5 correlation function across the Rindler-like horizon, satisfying the imaginary-time periodicity at $\beta = 2\pi/a$ from \S5, therefore satisfies the KMS condition at temperature
\[
T = a / (2\pi)
\]
This is the substrate-level Unruh temperature.

\subsection{Spectral content}

The KMS condition at temperature $T$ implies that the spectral function of the V5 correlations has a Bose-Einstein structure:
\[
n(\omega) = 1 / (e^{\omega/T} - 1) \quad \text{for } \omega > 0
\]
For substrate modes accessible to the accelerated observer in the Rindler-like wedge, the V5 vacuum looks like a \emph{thermal bath} at temperature $T = a/(2\pi)$. Each accessible mode has an occupation number given by the Planck distribution.

\subsection{What this delivers}

The substrate-level Unruh effect: a substrate observer accelerated relative to the asymptotic substrate frame at proper substrate-acceleration $a$ perceives the substrate's V5 vacuum as a thermal state at temperature $T = a/(2\pi)$. The thermalisation is not an artifact of the Bogoliubov calculation; it is the substrate-level structure that the analytic continuation of substrate-time imposes on V5 correlations in the accelerated frame.

The substrate observer is not in a thermal \emph{state}; the substrate is not heated. The thermalisation is \emph{frame-relative}: the asymptotic inertial substrate frame sees the V5 vacuum as the substrate-vacuum (no thermalisation), while the accelerated frame sees the same V5 vacuum as a thermal bath. The two views describe the same substrate-level reality through different observer-frame structures.

\section{DCGT Identification with the Standard Unruh Temperature}

The substrate calculation produces $T = a/(2\pi)$ where $a$ is the substrate-level proper acceleration. To recover the standard Unruh result $T_U = a/(2\pi)$, the substrate-level proper acceleration must identify with the standard proper acceleration.

\subsection{The DCGT identification}

DCGT (Diffusion Coarse-Graining Theorem) supplies the substrate-to-continuum bridge for substrate-state quantities. At leading-order coarse-graining:

\begin{itemize}[noitemsep]
\item Substrate-level proper substrate-acceleration $a_{\mathrm{substrate}} \to a$ (standard proper acceleration in the inertial frame).
\item Substrate-level proper substrate-time $t \to t$ (standard proper time).
\item V5 correlation function $\to$ standard Wightman correlation function in the appropriate limits.
\end{itemize}

The identification is FORCED at leading order via DCGT's substrate-to-continuum bridge applied to accelerated frames.

\subsection{The recovered Unruh temperature}

Combining the substrate calculation $T = a_{\mathrm{substrate}}/(2\pi)$ with the DCGT identification $a_{\mathrm{substrate}} = a$:
\[
T_U = a / (2\pi)
\]
This is the standard Unruh temperature. With explicit factors:
\[
T_U = \hbar a / (2\pi c k_B)
\]
For a hypothetical observer accelerated at $a = c^{2}/L$ (where $L$ is the proper distance to the Rindler horizon), the Unruh temperature is $T_U = \hbar c/(2\pi L k_B)$ --- inversely proportional to the proper distance to the horizon.

\subsection{Numerical scales}

For everyday accelerations the Unruh temperature is absurdly low:

\begin{center}
\begin{tabular}{ll}
\toprule
\textbf{Acceleration $a$} & \textbf{Unruh temperature $T_U$} \\
\midrule
$9.8$ m/s$^2$ (Earth surface gravity) & $\sim 4 \times 10^{-20}$ K \\
$10^{20}$ m/s$^2$ (large-particle accelerator) & $\sim 0.4$ K \\
$10^{26}$ m/s$^2$ (extreme laser fields) & $\sim 4 \times 10^{6}$ K \\
$a \sim c/\ell_P$ (Planck-scale acceleration) & $\sim 10^{32}$ K (Planck temperature) \\
\bottomrule
\end{tabular}
\end{center}

Direct detection of Unruh radiation has not been achieved; analog systems (BEC accelerating mirrors, accelerating photon detectors) have observed Unruh-analog signatures matching the predicted scaling.

\section{Substrate-Cutoff Corrections at High Frequencies}

The framework reproduces standard Unruh exactly at leading order. First-subleading-order corrections come from V5 finite-memory at the substrate-cutoff scale.

\subsection{The V5 cutoff}

The V5 kernel has finite memory time $\tau_{V5} = \ell_P/c$ at the gravitational scale (substrate-built timescale at the Planck scale). The kernel's frequency-domain magnitude:
\[
|\tilde V_5(\omega)|^{2} = (\mathcal{V}_0 \tau_{V5})^{2} / (1 + (\omega \tau_{V5})^{2})
\]
For $\omega \ll 1/\tau_{V5} = c/\ell_P$ (Planck frequency), V5 correlations are coherent. For $\omega \gtrsim c/\ell_P$, V5 coherence breaks down --- the substrate cannot mediate cross-chain correlations at frequencies approaching the Planck scale.

\subsection{The corrected Unruh spectrum}

The V5 cutoff modulates the Unruh spectrum at first subleading order:
\[
n_{\mathrm{ED}}(\omega) = n_U(\omega) / (1 + (\omega \tau_{V5})^{2})
\]
For Unruh frequencies $\omega \sim T_U \sim a$ (with $a$ in natural units): $\omega\tau_{V5} \sim a\ell_P/c$. For ordinary laboratory accelerations, $a\ell_P/c$ is many orders of magnitude smaller than unity, and the correction is invisible. Only at accelerations approaching $a \sim c^{2}/\ell_P$ --- Planck-scale accelerations --- does the V5 cutoff become significant.

\subsection{The trans-Unruh problem and its resolution}

The standard Unruh derivation has a structural problem analogous to the trans-Planckian problem of Hawking: modes observed at moderate frequencies in the accelerated frame trace back to arbitrarily-blueshifted modes near the Rindler-like horizon. At accelerations approaching the Planck scale, the standard derivation runs into structural difficulties.

The framework resolves this at the substrate level. V5 finite-memory cuts off mode coherence at $\omega \sim c/\ell_P$. The substrate cannot support arbitrarily-blueshifted modes near the Rindler-like horizon. The trans-Unruh problem is regulated naturally by V5's substrate-cutoff structure.

This is structurally identical to the V5 trans-Planckian resolution in the H-arc Hawking derivation. Both problems arise from the same Euclidean-periodicity argument applied at different horizon types; both are resolved by the same V5 substrate-cutoff mechanism.

\section{Cross-Domain Echoes}

The substrate-Unruh argument has structural echoes across multiple framework sectors.

\subsection{The Hawking echo}

The H-arc Hawking derivation runs the substrate-Unruh argument at the saturated decoupling surface (the substrate-level analog of a black-hole horizon). The substrate-level surface gravity $\kappa_{\mathrm{ED}} = \alpha(\nabla\sigma)|_{\mathrm{surf}}$ plays the role of the substrate-level proper acceleration. The same imaginary-time periodicity argument forces $\beta = 2\pi/\kappa_{\mathrm{ED}}$, giving $T_H = \kappa/(2\pi)$ via DCGT identification.

The two derivations are structurally identical --- same V5 vacuum, same imaginary-time periodicity argument, same KMS-condition equivalent. They differ only in \emph{which substrate horizon} the argument is applied to: the BH saturated decoupling surface (Hawking) or the accelerated-observer Rindler-like horizon (Unruh). The Hawking-Unruh equivalence in standard physics --- $T_H = \kappa/(2\pi)$ exactly parallels $T_U = a/(2\pi)$ --- emerges from this shared substrate-level mechanism.

\subsection{The de Sitter echo}

A static observer in de Sitter spacetime --- the spacetime of an accelerating universe with positive cosmological constant --- experiences a \emph{cosmological horizon} at proper distance $1/H$, where $H$ is the Hubble rate. The standard result (Gibbons-Hawking 1977) gives a thermal temperature
\[
T_{dS} = H / (2\pi)
\]
at the cosmological horizon.

The substrate-Unruh argument extends to this case directly. A substrate observer at fixed proper distance from the cosmological horizon experiences proper substrate-acceleration $H$, and the Rindler-like substrate frame near the horizon has the same structure as the BH and accelerated-observer cases. The imaginary-time periodicity argument forces $\beta = 2\pi/H$, giving $T_{dS} = H/(2\pi)$ via DCGT identification.

The substrate-Unruh walkthrough is therefore a structural prerequisite for any future ED arc on cosmological horizons / de Sitter thermality. The substrate mechanism is established here once and reused there.

\subsection{Analog Unruh experiments}

Analog-gravity experiments using BEC accelerating mirrors, accelerated photon detectors, and similar systems have been designed to probe Unruh-like thermalisation in laboratory-accessible regimes. The framework predicts:

\begin{itemize}[noitemsep]
\item Analog systems with proper acceleration $a_{\mathrm{analog}}$ and substrate-cutoff scale $\tau_{V5}^{\mathrm{analog}}$ (set by the analog system's microscopic correlation time) exhibit thermal spectra at $T = a_{\mathrm{analog}}/(2\pi)$.
\item The high-frequency tail is cut off at $\omega \sim 1/\tau_{V5}^{\mathrm{analog}}$ via the V5-form modulation $1/(1+(\omega\tau_{V5})^{2})$.
\item Cross-platform consistency: different analog systems should reproduce the same scaling structure with platform-specific cutoff scales.
\end{itemize}

These predictions parallel the analog-Hawking predictions of the H-arc and provide a falsifiable structural test of the framework's substrate-cutoff mechanism.

\subsection{The cross-domain unification}

V5 finite-memory does the substrate work in three apparently different contexts: Maxwell viscoelastic memory in soft matter (DCGT consequence), Hawking spectrum cutoff at black-hole horizons, and now Unruh radiation at accelerated-observer Rindler-like horizons. One substrate primitive, three physical phenomena. The cross-domain unification noted in the H-arc synthesis extends to include the substrate-Unruh effect as a fourth application.

\section{What's Forced, What's Inherited, What's Open}

It is worth being precise about what the framework's substrate-Unruh content delivers and what remains.

\subsection{What's forced}

The substrate-level Unruh temperature $T = a/(2\pi)$ at leading order is FORCED via the substrate-Unruh argument plus DCGT identification. The argument runs through:
\begin{itemize}[noitemsep]
\item Rindler-like substrate frame for accelerated observers (FORCED by substrate-level reading of acceleration plus T18 V1 forward-cone).
\item Imaginary-time periodicity at $\beta = 2\pi/a$ (FORCED by substrate-vacuum self-consistency at the Rindler-like horizon, parallel to no-conical-singularity argument).
\item KMS-condition equivalent (FORCED mathematical identity).
\item DCGT identification of substrate-level proper acceleration with standard proper acceleration (FORCED at leading order via DCGT).
\end{itemize}

The first-subleading-order corrections from V5 cutoff are FORM-FORCED via V5 finite-memory structure. The cutoff scale $\tau_{V5} = \ell_P/c$ at the gravitational scale is FORCED via T19 + dimensional analysis.

The structural identification of the Hawking and de Sitter cases as instances of the same substrate-Unruh argument is FORCED --- the three derivations share the same V5 substrate vacuum, the same imaginary-time periodicity argument, and the same KMS-condition equivalent. They differ only in which substrate horizon the argument applies to.

\subsection{What's inherited}

The numerical value of $\hbar$ in the explicit-units Unruh temperature ($T_U = \hbar a/(2\pi c k_B)$) is INHERITED from the dimensional atlas / Madelung anchoring.

The numerical value of $c$ is INHERITED from substrate constants.

The proper acceleration $a$ in any specific physical setup is INHERITED from the experimental context.

The cutoff scale $\tau_{V5}$ in non-gravitational contexts (analog systems) is INHERITED from the analog system's microscopic correlation time.

\subsection{What's open}

The closed-form derivation of $\tau_{V5}$ at the gravitational scale from substrate-microscopic V1 + V5 details is open. The framework identifies $\tau_{V5} = \ell_P/c$ via dimensional analysis on substrate primitives, but a closed-form derivation from more fundamental substrate timescales would tighten the cutoff identification.

The substrate-Unruh argument's extension to non-stationary accelerated frames (where the proper acceleration changes with proper time) is not yet developed in the framework. This is a downstream extension if needed for specific applications.

The empirical falsification of substrate-Unruh predictions awaits direct or analog Unruh detection. Current analog experiments are at the level of confirming the spectral form; precision tests of the substrate-cutoff modification require improved sensitivity.

The relationship between the substrate-Unruh argument and other thermalisation phenomena in the framework (entanglement entropy in Arc E, Hawking spectrum in Arc Hawking) is structurally clean but the detailed cross-arc unification has not been articulated as a publication-grade result yet.

\section{What This Argument Establishes}

The chain runs:

Substrate primitives (micro-events, chains, bandwidth, polarity, continuous time, locality) $\to$ V5 finite-memory kernel + V1 forward-cone-only (T18) supply substrate-level vacuum and causal structure $\to$ substrate-level reading of acceleration produces Rindler-like substrate frame for accelerated observer $\to$ analytic continuation $t \to -i\tau$ produces Euclidean substrate geometry with conical-point at Rindler-like horizon $\to$ no-substrate-conical-singularity argument forces $\beta = 2\pi/a$ $\to$ KMS condition gives V5 correlations thermal at $T = a/(2\pi)$ $\to$ DCGT identification recovers standard Unruh temperature $T_U = a/(2\pi)$ $\to$ first-subleading V5 cutoff at $\tau_{V5} = \ell_P/c$ regulates trans-Unruh problem.

The Unruh effect is now a derived consequence of substrate ontology rather than a postulate-level statement about quantum-field-theoretic Bogoliubov mixing. The mathematical content of the standard Unruh derivation --- the Rindler-frame analysis, the analytic continuation, the KMS condition, the thermal spectrum --- is unchanged. What changes is the foundational status: the Unruh temperature emerges from substrate-vacuum self-consistency at the Rindler-like horizon, not from a vacuum-state convention.

The framework reproduces standard Unruh exactly at observable scales. Laboratory acceleration regimes ($a \ll c^{2}/\ell_P$) match the standard Unruh prediction with V5 cutoff corrections invisible. Analog Unruh experiments confirm the spectral form. The substrate-cutoff resolution of the trans-Unruh problem is structurally analogous to the trans-Planckian resolution of standard Hawking: both are regulated naturally by V5's finite-memory substrate structure.

What's new is the substrate-level account of \emph{why} the Unruh thermalisation has the form it has. In standard physics, this is a Bogoliubov-transformation calculation with no underlying ontological account of why the Minkowski vacuum looks thermal to accelerated observers. In ED, it is the substrate-level statement that V5 cross-chain correlations near a Rindler-like horizon must be self-consistent under analytic continuation, and self-consistency forces imaginary-time periodicity at the inverse Unruh temperature.

The walkthrough is structurally identical to the H-1 piece of the Hawking arc. They are the same argument applied to different substrate horizons:
\begin{itemize}[noitemsep]
\item H-1 applies the substrate-Unruh argument at the saturated decoupling surface (BH horizon).
\item This walkthrough applies it at the accelerated-observer Rindler-like horizon (Unruh).
\end{itemize}

The substrate-Unruh argument is therefore a \emph{separable substrate-level result} that can be applied to any context with a participation-bandwidth boundary. Future arcs on cosmological horizons (de Sitter thermality), mesoscopic accelerated systems, or other accelerated-frame contexts inherit this walkthrough's substrate-Unruh content directly.

The factor that's worth emphasising: this walkthrough introduces no new substrate primitive. Every primitive used --- micro-events, chains, bandwidth, polarity, continuous time, locality, V1 forward-cone, V5 finite-memory --- was already in the framework's inventory from the QM-emergence and BH walkthroughs. The substrate-Unruh effect is what falls out when these primitives are applied to an accelerated frame. The substrate inventory is unchanged; the structural-foundations theorem inventory does not grow --- the Unruh effect is a downstream consequence of T18 + V5 + DCGT applied to the Rindler-like substrate frame, rather than a new theorem at the same structural level.

Whether the substrate primitives themselves are right is the load-bearing empirical question, as in every walkthrough. The framework stands or falls on whether participation, bandwidth, channels, polarity, locality, V1, V5, and the substrate-level structural commitments are the correct foundational concepts. The empirical exposure of the framework lives across closed sectors --- soft-matter mobility, substrate-derived gravity transitions, quantum-computational ceilings, Hawking radiation predictions --- not exclusively in the Unruh effect, where the framework reproduces standard QFT-in-flat-spacetime predictions plus FORM-FORCED first-subleading-order content.

For the Unruh effect specifically, the substrate-level case is closed at leading + first-subleading order. The framework reproduces standard Unruh at leading order via DCGT identification, supplies substrate-level UV regulation that resolves the trans-Unruh problem via V5 finite-memory, and produces FORM-FORCED first-subleading-order corrections distinguishing ED from strict semiclassical Unruh at extreme acceleration scales. Standard Unruh has been correct since 1976; ED supplies the substrate-level account of why.

\section{References}

\begin{itemize}[noitemsep]
\item Unruh, W. G. ``Notes on Black-Hole Evaporation.'' \textit{Physical Review D} \textbf{14}, 870--892 (1976).
\item Davies, P. C. W. ``Scalar Production in Schwarzschild and Rindler Metrics.'' \textit{Journal of Physics A} \textbf{8}, 609--616 (1975).
\item Fulling, S. A. ``Nonuniqueness of Canonical Field Quantization in Riemannian Space-Time.'' \textit{Physical Review D} \textbf{7}, 2850--2862 (1973).
\item Bisognano, J. J., Wichmann, E. H. ``On the Duality Condition for a Hermitian Scalar Field.'' \textit{Journal of Mathematical Physics} \textbf{16}, 985--1007 (1975).
\item Gibbons, G. W., Hawking, S. W. ``Cosmological Event Horizons, Thermodynamics, and Particle Creation.'' \textit{Physical Review D} \textbf{15}, 2738--2751 (1977).
\item Crispino, L. C. B., Higuchi, A., Matsas, G. E. A. ``The Unruh Effect and Its Applications.'' \textit{Reviews of Modern Physics} \textbf{80}, 787--838 (2008).
\item Kubo, R. ``Statistical-Mechanical Theory of Irreversible Processes I.'' \textit{Journal of the Physical Society of Japan} \textbf{12}, 570--586 (1957).
\item Martin, P. C., Schwinger, J. ``Theory of Many-Particle Systems I.'' \textit{Physical Review} \textbf{115}, 1342--1373 (1959).
\item Steinhauer, J. ``Observation of Quantum Hawking Radiation and Its Entanglement in an Analogue Black Hole.'' \textit{Nature Physics} \textbf{12}, 959--965 (2016).
\item Proxmire, A. \textit{Theorem 18: V1 Kernel Retardation and the Kernel-Level Arrow of Time.} April 2026.
\item Proxmire, A. \textit{Theorem 19: Newton's Law from Substrate Holographic Counting and the Identification of $\ell_P$.} April 2026.
\item Proxmire, A. \textit{The Diffusion Coarse-Graining Theorem: Substrate-to-Continuum Bridge for Canonical-ED Dynamical Content.} April 2026.
\item Proxmire, A. \textit{Arc Hawking H-1: Spectral Form and Temperature from V5 Cross-Chain Correlations.} May 2026.
\item Proxmire, A. \textit{Arc Hawking H-7: Synthesis and Cross-Domain Unification.} May 2026.
\item Proxmire, A. \textit{Walkthrough: From Primitives to Hawking Radiation.} May 2026.
\item Proxmire, A. \textit{Event Density: One Substrate, Three Domains.} April 2026.
\item The full Event Density corpus, including all forced theorems and supporting memos, is available at \url{https://github.com/allen-proxmire/event-density}.
\end{itemize}

\end{document}
